\chapter{Photopsia}

\section*{Definition}
Photopsia is the perception of flashes or sparkles of light in the visual field that are not caused by external light stimuli, often described as lightning streaks, camera flashes, or twinkling lights.

\section*{What Happens in Your Body}
Mechanical stimulation of the retina (particularly the peripheral retina) by traction, compression, or ischemia triggers depolarization of photoreceptors or retinal ganglion cells, which the brain interprets as light. Common mechanisms include vitreous traction on the retina during posterior vitreous detachment or direct neuronal irritation.

\section*{Causes}
\begin{itemize}
  \item \textbf{Posterior vitreous detachment (PVD):} Most common cause; vitreous gel separates from the retina, tugging at points of adhesion
  \item \textbf{Retinal tear or detachment:} Traction from PVD creates a full-thickness retinal break
  \item \textbf{Migraine aura:} Cortical spreading depression in occipital cortex (with or without headache)
  \item \textbf{Optic neuritis:} Demyelination causing conduction delay and phosphenes
  \item \textbf{Vertebrobasilar insufficiency:} Transient ischemia of the occipital lobe
  \item \textbf{Trauma:} Blunt ocular injury causing retinopathy (commotio retinae)
  \item \textbf{Toxic-metabolic:} Digoxin toxicity, LSD, or other hallucinogenic agents
\end{itemize}

\section*{When to See a Doctor}
See your doctor if:
\begin{itemize}
  \item Acute onset of photopsia, especially with a shower of floaters or curtain-like vision loss
  \item Photopsia persisting beyond 30 minutes or associated with headache
  \item History of high myopia, prior cataract surgery, or ocular trauma
  \item New-onset photopsia in an immunocompromised person (suspect infectious retinitis)
\end{itemize}

\section*{Related Symptoms}
\begin{itemize}
  \item Floaters (vitreous opacities)
  \item Visual field defect (scotoma)
  \item Blurred vision
  \item Headache (migraine-associated)
  \item Eye pain with movement (optic neuritis)
\end{itemize}
\textbf{Important:} Acute photopsia with new floaters is a retinal tear until proven otherwise; urgent dilated fundus examination with scleral depression is mandatory.

\section*{Self-Care / Home Management}
\begin{itemize}
  \item Avoid rapid head movements that may exacerbate vitreoretinal traction
  \item Dim the lights if photopsia is bothersome
  \item For known migraine aura: rest in a dark, quiet room until aura resolves
  \item If photopsia is recurrent and previously evaluated as benign, no specific home treatment
  \item Keep a symptom log (frequency, duration, triggers) to share with ophthalmologist
\end{itemize}

\section*{How Your Doctor Will Evaluate You}
\begin{itemize}
  \item Comprehensive history (laterality, character, duration, associated floaters or visual field loss)
  \item Dilated fundus examination with scleral depression of the peripheral retina
  \item Binocular indirect ophthalmoscopy to identify retinal tears, lattice degeneration, or PVD
  \item B-scan ultrasonography if vitreous hemorrhage precludes fundus visualization
  \item Optical coherence tomography (OCT) of the macula for vitreomacular traction
\end{itemize}

\section*{Treatment Options}
\begin{itemize}
  \item Posterior vitreous detachment without retinal tear: Observation with activity restriction and scheduled re-examination in 2--4 weeks
  \item Retinal tear: Laser retinopexy or cryotherapy to create chorioretinal adhesion around the break
  \item Retinal detachment: Immediate vitreoretinal surgical referral (pneumatic retinopexy, scleral buckle, or pars plana vitrectomy)
  \item Migraine-associated photopsia: Acute triptan therapy and prophylactic migraine management
  \item Optic neuritis: IV methylprednisolone followed by oral taper
\end{itemize}

\section*{References}
\begin{thebibliography}{99}
\bibitem{photopsia:ref1} Kahawita S, Flaherman M, Khan J. ``Posterior vitreous detachment and the risk of retinal tears.'' \textit{Journal of Clinical Medicine}, 2022;11(8):2131.
\bibitem{photopsia:ref2} Byer NE. ``Natural history of posterior vitreous detachment with early management as the premier line of defense against retinal detachment.'' \textit{Ophthalmology}, 1994;101(9):1503--1513.
\bibitem{photopsia:ref3} American Academy of Ophthalmology. \textit{Preferred Practice Pattern: Posterior Vitreous Detachment, Retinal Breaks, and Lattice Degeneration}. AAO, 2020.
\bibitem{photopsia:ref4} Gass JD, Norton EW. ``Photopsia and vitreous floaters.'' \textit{Archives of Ophthalmology}, 2000;118(7):997--1001.
\bibitem{photopsia:ref5} Russell MS, Gutteridge IF. ``Photopsia: a clinical review.'' \textit{Clinical and Experimental Optometry}, 2010;93(2):77--86.
\bibitem{photopsia:ref6} Singh P, Tripathy K. ``Photopsia: when should you worry?'' \textit{British Medical Journal}, 2021;374:n1785.

\end{thebibliography}
